\IfFileExists{stacks-project.cls}{%
\documentclass{stacks-project}
}{%
\documentclass{amsart}
}

% For dealing with references we use the comment environment
\usepackage{verbatim}
\newenvironment{reference}{\comment}{\endcomment}
%\newenvironment{reference}{}{}
\newenvironment{slogan}{\comment}{\endcomment}
\newenvironment{history}{\comment}{\endcomment}

% For commutative diagrams we use Xy-pic
\usepackage[all]{xy}

% We use 2cell for 2-commutative diagrams.
\xyoption{2cell}
\UseAllTwocells

% We use multicol for the list of chapters between chapters
\usepackage{multicol}

% This is generally recommended for better output
\usepackage{lmodern}
\usepackage[T1]{fontenc}

% For cross-file-references
\usepackage{xr-hyper}

% Package for hypertext links:
\usepackage{hyperref}

% For any local file, say "hello.tex" you want to link to please
% use \externaldocument[hello-]{hello}
\externaldocument[introduction-]{introduction}
\externaldocument[conventions-]{conventions}
\externaldocument[sets-]{sets}
\externaldocument[categories-]{categories}
\externaldocument[topology-]{topology}
\externaldocument[sheaves-]{sheaves}
\externaldocument[sites-]{sites}
\externaldocument[stacks-]{stacks}
\externaldocument[fields-]{fields}
\externaldocument[algebra-]{algebra}
\externaldocument[brauer-]{brauer}
\externaldocument[homology-]{homology}
\externaldocument[derived-]{derived}
\externaldocument[simplicial-]{simplicial}
\externaldocument[more-algebra-]{more-algebra}
\externaldocument[smoothing-]{smoothing}
\externaldocument[modules-]{modules}
\externaldocument[sites-modules-]{sites-modules}
\externaldocument[injectives-]{injectives}
\externaldocument[cohomology-]{cohomology}
\externaldocument[sites-cohomology-]{sites-cohomology}
\externaldocument[dga-]{dga}
\externaldocument[dpa-]{dpa}
\externaldocument[sdga-]{sdga}
\externaldocument[hypercovering-]{hypercovering}
\externaldocument[schemes-]{schemes}
\externaldocument[constructions-]{constructions}
\externaldocument[properties-]{properties}
\externaldocument[morphisms-]{morphisms}
\externaldocument[coherent-]{coherent}
\externaldocument[divisors-]{divisors}
\externaldocument[limits-]{limits}
\externaldocument[varieties-]{varieties}
\externaldocument[topologies-]{topologies}
\externaldocument[descent-]{descent}
\externaldocument[perfect-]{perfect}
\externaldocument[more-morphisms-]{more-morphisms}
\externaldocument[flat-]{flat}
\externaldocument[groupoids-]{groupoids}
\externaldocument[more-groupoids-]{more-groupoids}
\externaldocument[etale-]{etale}
\externaldocument[chow-]{chow}
\externaldocument[intersection-]{intersection}
\externaldocument[pic-]{pic}
\externaldocument[weil-]{weil}
\externaldocument[adequate-]{adequate}
\externaldocument[dualizing-]{dualizing}
\externaldocument[duality-]{duality}
\externaldocument[discriminant-]{discriminant}
\externaldocument[derham-]{derham}
\externaldocument[local-cohomology-]{local-cohomology}
\externaldocument[algebraization-]{algebraization}
\externaldocument[curves-]{curves}
\externaldocument[resolve-]{resolve}
\externaldocument[models-]{models}
\externaldocument[functors-]{functors}
\externaldocument[equiv-]{equiv}
\externaldocument[pione-]{pione}
\externaldocument[etale-cohomology-]{etale-cohomology}
\externaldocument[proetale-]{proetale}
\externaldocument[relative-cycles-]{relative-cycles}
\externaldocument[more-etale-]{more-etale}
\externaldocument[trace-]{trace}
\externaldocument[crystalline-]{crystalline}
\externaldocument[spaces-]{spaces}
\externaldocument[spaces-properties-]{spaces-properties}
\externaldocument[spaces-morphisms-]{spaces-morphisms}
\externaldocument[decent-spaces-]{decent-spaces}
\externaldocument[spaces-cohomology-]{spaces-cohomology}
\externaldocument[spaces-limits-]{spaces-limits}
\externaldocument[spaces-divisors-]{spaces-divisors}
\externaldocument[spaces-over-fields-]{spaces-over-fields}
\externaldocument[spaces-topologies-]{spaces-topologies}
\externaldocument[spaces-descent-]{spaces-descent}
\externaldocument[spaces-perfect-]{spaces-perfect}
\externaldocument[spaces-more-morphisms-]{spaces-more-morphisms}
\externaldocument[spaces-flat-]{spaces-flat}
\externaldocument[spaces-groupoids-]{spaces-groupoids}
\externaldocument[spaces-more-groupoids-]{spaces-more-groupoids}
\externaldocument[bootstrap-]{bootstrap}
\externaldocument[spaces-pushouts-]{spaces-pushouts}
\externaldocument[spaces-chow-]{spaces-chow}
\externaldocument[groupoids-quotients-]{groupoids-quotients}
\externaldocument[spaces-more-cohomology-]{spaces-more-cohomology}
\externaldocument[spaces-simplicial-]{spaces-simplicial}
\externaldocument[spaces-duality-]{spaces-duality}
\externaldocument[formal-spaces-]{formal-spaces}
\externaldocument[restricted-]{restricted}
\externaldocument[spaces-resolve-]{spaces-resolve}
\externaldocument[formal-defos-]{formal-defos}
\externaldocument[defos-]{defos}
\externaldocument[cotangent-]{cotangent}
\externaldocument[examples-defos-]{examples-defos}
\externaldocument[algebraic-]{algebraic}
\externaldocument[examples-stacks-]{examples-stacks}
\externaldocument[stacks-sheaves-]{stacks-sheaves}
\externaldocument[criteria-]{criteria}
\externaldocument[artin-]{artin}
\externaldocument[quot-]{quot}
\externaldocument[stacks-properties-]{stacks-properties}
\externaldocument[stacks-morphisms-]{stacks-morphisms}
\externaldocument[stacks-limits-]{stacks-limits}
\externaldocument[stacks-cohomology-]{stacks-cohomology}
\externaldocument[stacks-perfect-]{stacks-perfect}
\externaldocument[stacks-introduction-]{stacks-introduction}
\externaldocument[stacks-more-morphisms-]{stacks-more-morphisms}
\externaldocument[stacks-geometry-]{stacks-geometry}
\externaldocument[moduli-]{moduli}
\externaldocument[moduli-curves-]{moduli-curves}
\externaldocument[examples-]{examples}
\externaldocument[exercises-]{exercises}
\externaldocument[guide-]{guide}
\externaldocument[desirables-]{desirables}
\externaldocument[coding-]{coding}
\externaldocument[obsolete-]{obsolete}
\externaldocument[fdl-]{fdl}
\externaldocument[index-]{index}

% Theorem environments.
%
\theoremstyle{plain}
\newtheorem{theorem}[subsection]{Theorem}
\newtheorem{proposition}[subsection]{Proposition}
\newtheorem{lemma}[subsection]{Lemma}

\theoremstyle{definition}
\newtheorem{definition}[subsection]{Definition}
\newtheorem{example}[subsection]{Example}
\newtheorem{exercise}[subsection]{Exercise}
\newtheorem{situation}[subsection]{Situation}

\theoremstyle{remark}
\newtheorem{remark}[subsection]{Remark}
\newtheorem{remarks}[subsection]{Remarks}

\numberwithin{equation}{subsection}

% Macros
%
\def\lim{\mathop{\mathrm{lim}}\nolimits}
\def\colim{\mathop{\mathrm{colim}}\nolimits}
\def\Spec{\mathop{\mathrm{Spec}}}
\def\Hom{\mathop{\mathrm{Hom}}\nolimits}
\def\Ext{\mathop{\mathrm{Ext}}\nolimits}
\def\SheafHom{\mathop{\mathcal{H}\!\mathit{om}}\nolimits}
\def\SheafExt{\mathop{\mathcal{E}\!\mathit{xt}}\nolimits}
\def\Sch{\mathit{Sch}}
\def\Mor{\mathop{\mathrm{Mor}}\nolimits}
\def\Ob{\mathop{\mathrm{Ob}}\nolimits}
\def\Sh{\mathop{\mathit{Sh}}\nolimits}
\def\NL{\mathop{N\!L}\nolimits}
\def\CH{\mathop{\mathrm{CH}}\nolimits}
\def\proetale{{pro\text{-}\acute{e}tale}}
\def\etale{{\acute{e}tale}}
\def\QCoh{\mathit{QCoh}}
\def\Ker{\mathop{\mathrm{Ker}}}
\def\Im{\mathop{\mathrm{Im}}}
\def\Coker{\mathop{\mathrm{Coker}}}
\def\Coim{\mathop{\mathrm{Coim}}}

% Boxtimes
%
\DeclareMathSymbol{\boxtimes}{\mathbin}{AMSa}{"02}

%
% Macros for moduli stacks/spaces
%
\def\QCohstack{\mathcal{QC}\!\mathit{oh}}
\def\Cohstack{\mathcal{C}\!\mathit{oh}}
\def\Spacesstack{\mathcal{S}\!\mathit{paces}}
\def\Quotfunctor{\mathrm{Quot}}
\def\Hilbfunctor{\mathrm{Hilb}}
\def\Curvesstack{\mathcal{C}\!\mathit{urves}}
\def\Polarizedstack{\mathcal{P}\!\mathit{olarized}}
\def\Complexesstack{\mathcal{C}\!\mathit{omplexes}}
% \Pic is the operator that assigns to X its picard group, usage \Pic(X)
% \Picardstack_{X/B} denotes the Picard stack of X over B
% \Picardfunctor_{X/B} denotes the Picard functor of X over B
\def\Pic{\mathop{\mathrm{Pic}}\nolimits}
\def\Picardstack{\mathcal{P}\!\mathit{ic}}
\def\Picardfunctor{\mathrm{Pic}}
\def\Deformationcategory{\mathcal{D}\!\mathit{ef}}
\usepackage{cleveref}
\crefname{originaldefinition}{original definition}{original definitions}
\Crefname{originaldefinition}{Original definition}{Original definitions}
\makeatletter
\newlabel{morphisms-definition-finite-type-point}{{02J1}{source}{}{sourcepointer-1}{}}
\newlabel{morphisms-definition-finite-type-point@cref}{{[originaldefinition][0][]02J1}{[1][0][]source}{}{}{}}
\makeatother
\crefname{originalreference}{original reference}{original references}
\Crefname{originalreference}{Original reference}{Original references}
\makeatletter
\newlabel{situation-slice}{{04MY}{source}{}{sourcepointer-2}{}}
\newlabel{situation-slice@cref}{{[originalreference][0][]04MY}{[1][0][]source}{}{}{}}
\makeatother
\crefname{originallemma}{original lemma}{original lemmas}
\Crefname{originallemma}{Original lemma}{Original lemmas}
\makeatletter
\newlabel{lemma-slice}{{0461}{source}{}{sourcepointer-3}{}}
\newlabel{lemma-slice@cref}{{[originallemma][0][]0461}{[1][0][]source}{}{}{}}
\makeatother
\crefname{originalremark}{original remark}{original remarks}
\Crefname{originalremark}{Original remark}{Original remarks}
\makeatletter
\newlabel{remark-local-source-apply}{{047W}{source}{}{sourcepointer-4}{}}
\newlabel{remark-local-source-apply@cref}{{[originalremark][0][]047W}{[1][0][]source}{}{}{}}
\makeatother
\crefname{originallemma}{original lemma}{original lemmas}
\Crefname{originallemma}{Original lemma}{Original lemmas}
\makeatletter
\newlabel{lemma-restrict-property}{{04MV}{source}{}{sourcepointer-5}{}}
\newlabel{lemma-restrict-property@cref}{{[originallemma][0][]04MV}{[1][0][]source}{}{}{}}
\makeatother
\crefname{originallemma}{original lemma}{original lemmas}
\Crefname{originallemma}{Original lemma}{Original lemmas}
\makeatletter
\newlabel{morphisms-lemma-base-change-finite-presentation}{{01TS}{source}{}{sourcepointer-6}{}}
\newlabel{morphisms-lemma-base-change-finite-presentation@cref}{{[originallemma][0][]01TS}{[1][0][]source}{}{}{}}
\makeatother
\crefname{originallemma}{original lemma}{original lemmas}
\Crefname{originallemma}{Original lemma}{Original lemmas}
\makeatletter
\newlabel{more-morphisms-lemma-base-change-CM}{{045T}{source}{}{sourcepointer-7}{}}
\newlabel{more-morphisms-lemma-base-change-CM@cref}{{[originallemma][0][]045T}{[1][0][]source}{}{}{}}
\makeatother
\crefname{originallemma}{original lemma}{original lemmas}
\Crefname{originallemma}{Original lemma}{Original lemmas}
\makeatletter
\newlabel{morphisms-lemma-identify-finite-type-points}{{02J2}{source}{}{sourcepointer-8}{}}
\newlabel{morphisms-lemma-identify-finite-type-points@cref}{{[originallemma][0][]02J2}{[1][0][]source}{}{}{}}
\makeatother
\crefname{originallemma}{original lemma}{original lemmas}
\Crefname{originallemma}{Original lemma}{Original lemmas}
\makeatletter
\newlabel{lemma-double-restrict-property}{{04MX}{source}{}{sourcepointer-9}{}}
\newlabel{lemma-double-restrict-property@cref}{{[originallemma][0][]04MX}{[1][0][]source}{}{}{}}
\makeatother
\crefname{originallemma}{original lemma}{original lemmas}
\Crefname{originallemma}{Original lemma}{Original lemmas}
\makeatletter
\newlabel{more-morphisms-lemma-CM-local-source-and-target}{{045V}{source}{}{sourcepointer-10}{}}
\newlabel{more-morphisms-lemma-CM-local-source-and-target@cref}{{[originallemma][0][]045V}{[1][0][]source}{}{}{}}
\makeatother
\title{Maximal slicing of a Cohen-Macaulay groupoid scheme at a finite type point, keeping h, s' and t' Cohen-Macaulay and locally of finite presentation}
\author{The Stacks Project Authors}
\date{stacks-project:more-groupoids.tex:local-snapshot-20261012; GNU FDL 1.2 or later}
\begin{document}
\maketitle
\setcounter{subsection}{1}
\noindent\textit{Editorial note.} All delivered mathematical text is the author's own, copied from the pinned Stacks Project chapter file more-groupoids.tex; no statement or proof step was rewritten, reordered or removed. The counted span is the author environment and its own proof environment. References inside the span are handed to the assembler as navigation stubs that carry the canonical tag of the cited tag only; whether each of them is load-bearing is recorded on the stub itself.\par\medskip
\noindent\textit{Allowed standard prerequisites (editorial).}\begin{itemize}
\item Definition of a finite type point of a scheme (morphisms chapter, Definition morphisms-definition-finite-type-point, canonical tag 02J1). cited at statement lines 2251-2252. Fixes the class of points u of U at which the slicing problem is posed.
\item The slicing situation: notation F\_u, G\_u, G the stabilizer group scheme, and the objects F'\_u, G'\_u, R', s', t', h of a restriction (more-groupoids chapter, Situation situation-slice, canonical tag 04MY; kind reference). cited at statement lines 2253-2254. Supplies the notation in which the conclusion is phrased: the fibres and their dimensions, the stabilizer dimension dim(G'\_u), and the morphisms s', t' and h of the restricted groupoid.
\item The slice lemma: one slicing step at u giving an immersion g with properties (1), (2), (3), a strict orbit dimension drop, and h Cohen-Macaulay at (u, e(u)) (more-groupoids chapter, Lemma lemma-slice, canonical tag 0461). cited at proof lines 2277-2285. The workhorse of the induction step: produces the single slice whose h is Cohen-Macaulay at (u, e(u)) and whose fibre dimension strictly drops.
\item Remark producing the largest open on which a morphism locally of finite presentation is Cohen-Macaulay (more-groupoids chapter, Remark remark-local-source-apply, canonical tag 047W). cited at proof lines 2285-2290. Turns Cohen-Macaulayness of h at the single point (u, e(u)) into Cohen-Macaulayness of h on all of U'', by shrinking to the largest such open, which contains u.
\item Restriction lemma transferring properties of h to s' and t' of the restricted groupoid (more-groupoids chapter, Lemma lemma-restrict-property, canonical tag 04MV). cited at proof lines 2293-2295. Concludes that s' and t' are locally of finite presentation and Cohen-Macaulay once h is, via the two base change lemmas.
\item Base change of morphisms locally of finite presentation (morphisms chapter, Lemma morphisms-lemma-base-change-finite-presentation, canonical tag 01TS). cited at proof line 2296. Input to the restrict-property application: local finite presentation of s and t passes to the fibre products over which s' and t' factor.
\item Base change of Cohen-Macaulay morphisms (more-morphisms chapter, Lemma more-morphisms-lemma-base-change-CM, canonical tag 045T). cited at proof lines 2297-2298. Input to the restrict-property application: Cohen-Macaulayness of h passes to the base changes over which s' and t' factor.
\item Identification of finite type points under a locally finite presentation morphism (morphisms chapter, Lemma morphisms-lemma-identify-finite-type-points, canonical tag 02J2). cited at proof lines 2303-2306. Shows u remains a finite type point of U' after the immersion g, so the induction hypothesis applies to the restricted groupoid at u.
\item Double restriction lemma: the composition h'' for a composed slicing stays locally of finite presentation and Cohen-Macaulay (more-groupoids chapter, Lemma lemma-double-restrict-property, canonical tag 04MX). cited at proof lines 2318-2319. Verifies property (4) for the composed solution g'' = g composed with g'.
\item Cohen-Macaulay morphisms are fppf local on the target (more-morphisms chapter, Lemma more-morphisms-lemma-CM-local-source-and-target, canonical tag 045V). cited at proof lines 2321-2322. Used inside the double restrict-property application to descend Cohen-Macaulayness along the fppf covering of the target.
\end{itemize}\medskip
\ifdefined\hypertarget\hypertarget{sourcepointer-1}{}\fi
\noindent\textit{Source reference.} Original definition 02J1: sources/stacks-chow/tags.txt, source lines 3320--3320. Cited at statement lines 2251-2252: u is taken to be a finite type point of U, the class of points at which the slicing problem is posed.\par\medskip
\ifdefined\hypertarget\hypertarget{sourcepointer-2}{}\fi
\noindent\textit{Source reference.} Original reference 04MY: sources/stacks-chow/tags.txt, source lines 6059--6059. Cited at statement lines 2253-2254: notation as in Situation slice supplies F\_u, G\_u, the stabilizer group scheme G, and the restriction objects F'\_u, G'\_u, R', s', t' in which the six listed properties are phrased.\par\medskip
\ifdefined\hypertarget\hypertarget{sourcepointer-3}{}\fi
\noindent\textit{Source reference.} Original lemma 0461: sources/stacks-chow/tags.txt, source lines 5456--5456. Cited at proof lines 2277-2285: the induction step slices once at u, getting the immersion g with properties (1), (2), (3), the strict orbit dimension drop, and h Cohen-Macaulay at (u, e(u)).\par\medskip
\ifdefined\hypertarget\hypertarget{sourcepointer-4}{}\fi
\noindent\textit{Source reference.} Original remark 047W: sources/stacks-chow/tags.txt, source lines 5523--5523. Cited at proof lines 2285-2290: produces the largest open U'' of U' on which h is Cohen-Macaulay, containing (u, e(u)) and hence u, after which h is Cohen-Macaulay everywhere.\par\medskip
\ifdefined\hypertarget\hypertarget{sourcepointer-5}{}\fi
\noindent\textit{Source reference.} Original lemma 04MV: sources/stacks-chow/tags.txt, source lines 6056--6056. Cited at proof lines 2293-2295: transfers local finite presentation and Cohen-Macaulayness from h to s' and t' of the restriction.\par\medskip
\ifdefined\hypertarget\hypertarget{sourcepointer-6}{}\fi
\noindent\textit{Source reference.} Original lemma 01TS: sources/stacks-chow/tags.txt, source lines 2385--2385. Cited at proof line 2296 inside the restrict-property application: local finite presentation passes along base change.\par\medskip
\ifdefined\hypertarget\hypertarget{sourcepointer-7}{}\fi
\noindent\textit{Source reference.} Original lemma 045T: sources/stacks-chow/tags.txt, source lines 5448--5448. Cited at proof lines 2297-2298 inside the restrict-property application: Cohen-Macaulayness passes along base change.\par\medskip
\ifdefined\hypertarget\hypertarget{sourcepointer-8}{}\fi
\noindent\textit{Source reference.} Original lemma 02J2: sources/stacks-chow/tags.txt, source lines 3321--3321. Cited at proof lines 2303-2306: u is still a finite type point of U' after the immersion g, so the induction hypothesis applies at u.\par\medskip
\ifdefined\hypertarget\hypertarget{sourcepointer-9}{}\fi
\noindent\textit{Source reference.} Original lemma 04MX: sources/stacks-chow/tags.txt, source lines 6058--6058. Cited at proof lines 2318-2319: h'' for the composed slicing is locally of finite presentation and Cohen-Macaulay, giving property (4) of the composed solution.\par\medskip
\ifdefined\hypertarget\hypertarget{sourcepointer-10}{}\fi
\noindent\textit{Source reference.} Original lemma 045V: sources/stacks-chow/tags.txt, source lines 5450--5450. Cited at proof lines 2321-2322 inside the double restrict-property application: Cohen-Macaulay morphisms are fppf local on the target.\par\medskip
\begin{lemma}
\label{lemma-max-slice}
Let $S$ be a scheme.
Let $(U, R, s, t, c, e, i)$ be a groupoid scheme over $S$.
Let $G \to U$ be the stabilizer group scheme.
Assume $s$ and $t$ are Cohen-Macaulay and locally of finite presentation.
Let $u \in U$ be a finite type point of the scheme $U$, see
Morphisms, Definition \ref{morphisms-definition-finite-type-point}.
With notation as in
Situation \ref{situation-slice}
there exist an affine scheme $U'$ and a morphism $g : U' \to U$ such that
\begin{enumerate}
\item $g$ is an immersion,
\item $u \in U'$,
\item $g$ is locally of finite presentation,
\item the morphism $h : U' \times_{g, U, t} R \longrightarrow U$
is Cohen-Macaulay and locally of finite presentation,
\item the morphisms $s', t' : R' \to U'$ are Cohen-Macaulay and
locally of finite presentation, and
\item $\dim_{e(u)}(F'_u) = \dim(G'_u)$.
\end{enumerate}
\end{lemma}

\begin{proof}
As $s$ is locally of finite presentation the scheme $F_u$ is
locally of finite type over $\kappa(u)$. Hence
$\dim_{e(u)}(F_u) < \infty$ and we may argue by induction on
$\dim_{e(u)}(F_u)$.

\medskip\noindent
If $\dim_{e(u)}(F_u) = \dim(G_u)$ there is nothing to prove.
Assume $\dim_{e(u)}(F_u) > \dim(G_u)$. This means that
Lemma \ref{lemma-slice}
applies and we find a morphism $g : U' \to U$ which has
properties (1), (2), (3), instead of (6) we have
$\dim_{e(u)}(F'_u) < \dim_{e(u)}(F_u)$,
and instead of (4) and (5) we have that the composition
$$
h = s \circ \text{pr}_1 : U' \times_{g, U, t} R \longrightarrow U
$$
is Cohen-Macaulay at the point $(u, e(u))$. We apply
Remark \ref{remark-local-source-apply}
and we obtain an open subscheme $U'' \subset U'$ such that
$U'' \times_{g, U, t} R \subset U' \times_{g, U, t} R$
is the largest open subscheme on which $h$ is Cohen-Macaulay.
Since $(u, e(u)) \in U'' \times_{g, U, t} R$ we see that $u \in U''$.
Hence we may replace $U'$ by $U''$ and assume that in fact $h$ is
Cohen-Macaulay everywhere! By
Lemma \ref{lemma-restrict-property}
we conclude that $s', t'$ are locally of finite
presentation and Cohen-Macaulay (use
Morphisms, Lemma \ref{morphisms-lemma-base-change-finite-presentation}
and
More on Morphisms, Lemma \ref{more-morphisms-lemma-base-change-CM}).

\medskip\noindent
By construction $\dim_{e'(u)}(F'_u) < \dim_{e(u)}(F_u)$,
so we may apply the induction hypothesis to $(U', R', s', t', c')$
and the point $u \in U'$. Note that $u$ is also a finite type point
of $U'$ (for example you can see this using the characterization of
finite type points from
Morphisms, Lemma \ref{morphisms-lemma-identify-finite-type-points}).
Let $g' : U'' \to U'$ and $(U'', R'', s'', t'', c'')$ be the solution
of the corresponding problem starting with $(U', R', s', t', c')$
and the point $u \in U'$. We claim that the composition
$$
g'' = g \circ g' : U'' \longrightarrow U
$$
is a solution for the original problem. Properties (1), (2), (3), (5),
and (6) are immediate. To see (4) note that the morphism
$$
h'' = s \circ \text{pr}_1 : U'' \times_{g'', U, t} R \longrightarrow U
$$
is locally of finite presentation and Cohen-Macaulay by an application of
Lemma \ref{lemma-double-restrict-property}
(use
More on Morphisms, Lemma \ref{more-morphisms-lemma-CM-local-source-and-target}
to see that Cohen-Macaulay morphisms are fppf local on the target).
\end{proof}

\end{document}
